\section{Conclusiones}

\textbf{Big Data no es un tamaño, es una relación.} El mismo archivo de 38.2 millones de filas fue imposible en una máquina de 15.7 GB de RAM y manejable en una de 8 GB. Ninguna cifra absoluta de filas o gigabytes marca la frontera; la marca la interacción entre el volumen, la forma en que están escritos los datos y el equipo que los procesa.

\textbf{La representación es la palanca más barata.} El paso de CSV a Parquet con proyección y tipos reducidos costó entre 2.6 y 14.8 segundos por versión y multiplicó el rendimiento por factores de hasta $110\times$. Más importante: movió el umbral de viabilidad. Un equipo que no podía procesar la versión completa pasó a hacerlo en 6 segundos, sin agregar hardware.

\textbf{La arquitectura del motor importa más que sus especificaciones.} DuckDB, columnar y vectorizado, resolvió en 2.7 segundos sobre 8 GB de RAM lo que SQLite, orientado a filas, no terminaba en 10 minutos sobre 40 GB. Los índices y los \texttt{PRAGMA} mejoran a SQLite, pero no cambian su modelo de almacenamiento: un B-tree no convierte un motor OLTP en columnar.

\textbf{La carga, no la consulta, es el cuello de botella.} Entre el 57\% y el 97\% del tiempo total se va en leer el archivo. Cualquier esfuerzo de optimización que ignore esto ataca el problema equivocado.

\textbf{Sobre el sistema de bicicletas.} El uso es bimodal en días laborables, con picos a las 08:00 y las 17:00-18:00, y unimodal y recreativo en fin de semana, con máximo a las 14:00. El desbalance logístico sigue un patrón radial: las estaciones de tren y las zonas residenciales se vacían, los distritos de oficinas se saturan. Y el tráfico está concentrado pero no dominado: las 1,000 rutas más usadas suman el 9.6\% de los viajes, y las nueve primeras son circulares dentro de parques, es decir, paseos y no desplazamientos.
